\documentclass[10pt,a4paper]{amsart}
\usepackage{amsmath,amssymb,amsthm,graphicx}
\usepackage{verbatim}
\newenvironment{reference}{\comment}{\endcomment}
%\newenvironment{reference}{}{}
\newenvironment{slogan}{\comment}{\endcomment}
\newenvironment{history}{\comment}{\endcomment}

% For commutative diagrams we use Xy-pic
\usepackage[all]{xy}

% We use 2cell for 2-commutative diagrams.
\xyoption{2cell}
% Theorem environments.
%
\theoremstyle{plain}
\newtheorem{theorem}[subsection]{Theorem}
\newtheorem{proposition}[subsection]{Proposition}
\newtheorem{lemma}[subsection]{Lemma}

\theoremstyle{definition}
\newtheorem{definition}[subsection]{Definition}
\newtheorem{example}[subsection]{Example}
\newtheorem{exercise}[subsection]{Exercise}
\newtheorem{situation}[subsection]{Situation}

\theoremstyle{remark}
\newtheorem{remark}[subsection]{Remark}
\newtheorem{remarks}[subsection]{Remarks}

\numberwithin{equation}{subsection}

% Macros
%
\def\lim{\mathop{\mathrm{lim}}\nolimits}
\def\colim{\mathop{\mathrm{colim}}\nolimits}
\def\Spec{\mathop{\mathrm{Spec}}}
\def\Hom{\mathop{\mathrm{Hom}}\nolimits}
\def\Ext{\mathop{\mathrm{Ext}}\nolimits}
\def\SheafHom{\mathop{\mathcal{H}\!\mathit{om}}\nolimits}
\def\SheafExt{\mathop{\mathcal{E}\!\mathit{xt}}\nolimits}
\def\Sch{\mathit{Sch}}
\def\Mor{\mathop{\mathrm{Mor}}\nolimits}
\def\Ob{\mathop{\mathrm{Ob}}\nolimits}
\def\Sh{\mathop{\mathit{Sh}}\nolimits}
\def\NL{\mathop{N\!L}\nolimits}
\def\CH{\mathop{\mathrm{CH}}\nolimits}
\def\proetale{{pro\text{-}\acute{e}tale}}
\def\etale{{\acute{e}tale}}
\def\QCoh{\mathit{QCoh}}
\def\Ker{\mathop{\mathrm{Ker}}}
\def\Im{\mathop{\mathrm{Im}}}
\def\Coker{\mathop{\mathrm{Coker}}}
\def\Coim{\mathop{\mathrm{Coim}}}

% Boxtimes
%
\DeclareMathSymbol{\boxtimes}{\mathbin}{AMSa}{"02}

%
% Macros for moduli stacks/spaces
%
\def\QCohstack{\mathcal{QC}\!\mathit{oh}}
\def\Cohstack{\mathcal{C}\!\mathit{oh}}
\def\Spacesstack{\mathcal{S}\!\mathit{paces}}
\def\Quotfunctor{\mathrm{Quot}}
\def\Hilbfunctor{\mathrm{Hilb}}
\def\Curvesstack{\mathcal{C}\!\mathit{urves}}
\def\Polarizedstack{\mathcal{P}\!\mathit{olarized}}
\def\Complexesstack{\mathcal{C}\!\mathit{omplexes}}
% \Pic is the operator that assigns to X its picard group, usage \Pic(X)
% \Picardstack_{X/B} denotes the Picard stack of X over B
% \Picardfunctor_{X/B} denotes the Picard functor of X over B
\def\Pic{\mathop{\mathrm{Pic}}\nolimits}
\def\Picardstack{\mathcal{P}\!\mathit{ic}}
\def\Picardfunctor{\mathrm{Pic}}
\def\Deformationcategory{\mathcal{D}\!\mathit{ef}}
\usepackage{cleveref}
\crefname{originallemma}{original lemma}{original lemmas}
\Crefname{originallemma}{Original lemma}{Original lemmas}
\makeatletter
\newlabel{constructions-lemma-globally-generated-omega-twist-1}{{0B3B}{source}{}{sourcepointer-1}{}}
\newlabel{constructions-lemma-globally-generated-omega-twist-1@cref}{{[originallemma][0][]0B3B}{[1][0][]source}{}{}{}}
\makeatother
\crefname{originalreference}{original reference}{original references}
\Crefname{originalreference}{Original reference}{Original references}
\makeatletter
\newlabel{lemma-globally-generated}{{01AN}{source}{}{sourcepointer-2}{}}
\newlabel{lemma-globally-generated@cref}{{[originalreference][0][]01AN}{[1][0][]source}{}{}{}}
\makeatother
\crefname{originalreference}{original reference}{original references}
\Crefname{originalreference}{Original reference}{Original references}
\makeatletter
\newlabel{situation-rational}{{0B4W}{source}{}{sourcepointer-3}{}}
\newlabel{situation-rational@cref}{{[originalreference][0][]0B4W}{[1][0][]source}{}{}{}}
\makeatother
\crefname{originalreference}{original reference}{original references}
\Crefname{originalreference}{Original reference}{Original references}
\makeatletter
\newlabel{situation-vanishing}{{07UY}{source}{}{sourcepointer-4}{}}
\newlabel{situation-vanishing@cref}{{[originalreference][0][]07UY}{[1][0][]source}{}{}{}}
\makeatother
\title{Sections and vanishing of the first cohomology of powers of the ideal sheaf of the exceptional divisor}
\author{The Stacks Project Authors}
\date{stacks-project:resolve.tex:local-snapshot-20261012; GNU FDL 1.2 or later}
\begin{document}
\maketitle
\noindent\textit{Editorial note.} All delivered mathematical text is the author's own, copied from the pinned Stacks Project chapter file resolve.tex; no statement or proof step was rewritten, reordered or removed.\par\medskip
\noindent\textit{Editorial note.} The counted proof is the author's complete argument: the identification of the sections of the ideal sheaf with the maximal ideal, the vanishing through global generation, and the induction with the twisted kernel of the presentation.\par\medskip
\noindent\textit{Editorial note.} The label lemma-sections-powers-I-rational names the counted statement, and the proof environment immediately following it in the source is delivered verbatim as its proof.\par\medskip
\noindent\textit{Editorial note.} The counted statement is carried by the author's own lemma environment in the chapter source.\par\medskip
\noindent\textit{Editorial note.} Pre-open gate: python operations/agent-authoring/tools/precard\_gate.py --source resolve.tex --label lemma-sections-powers-I-rational --start 2487 --end 2531 --local-id 28099 returned RESULT=OK (L28099 free, occupancy FREE, public interval clear). Pre-assembly tag self-check: canonical tag resolved and not present in reference/public-items.jsonl.\par\medskip
\noindent\textit{Editorial note.} The statement and proof are copied from the pinned chapter file; the cited inputs are kept as external pointers with their canonical Stacks tags.\par\medskip
\noindent\textit{Allowed standard prerequisites (editorial).}\begin{itemize}
\item Globally generated modules have vanishing first cohomology in the rational situation. Same chapter, the lemma that a globally generated module on the resolution has vanishing first cohomology. Supplies the vanishing of the first cohomology of the powers of the ideal sheaf and of the auxiliary subsheaf used in the induction.
\item The ideal sheaf of the exceptional divisor is generated by the generators of the maximal ideal. The description of the blowup in terms of the Rees algebra, so that the powers of the ideal sheaf are the powers of the maximal ideal extended to the resolution. Identifies the powers of the ideal sheaf with the extended powers of the maximal ideal.
\item The tensor product of the auxiliary sheaf with the ideal sheaf is globally generated. Constructions, the lemma computing the twist of the kernel of the evaluation map on the resolution. Provides global generation of the tensor powers needed to run the induction.
\end{itemize}\medskip
\ifdefined\hypertarget\hypertarget{sourcepointer-1}{}\fi
\noindent\textit{Source reference.} Original lemma 0B3B: sources/stacks-chow/tags.txt, source lines 13959--13959. Prerequisite pointer invoked at proof line(s) 2515; it names a standard fact of a sibling chapter that the counted argument uses, and every step of the counted proof built on it is retained verbatim in the delivered body.\par\medskip
\ifdefined\hypertarget\hypertarget{sourcepointer-2}{}\fi
\noindent\textit{Source reference.} Original reference 01AN: sources/stacks-resolve/pinned-original-source.tex, source lines 2465--2465. Prerequisite pointer invoked at proof line(s) 2503; it names a standard fact of the same chapter that the counted argument uses, and every step of the counted proof built on it is retained verbatim in the delivered body.\par\medskip
\ifdefined\hypertarget\hypertarget{sourcepointer-3}{}\fi
\noindent\textit{Source reference.} Original reference 0B4W: sources/stacks-resolve/pinned-original-source.tex, source lines 2455--2455. Prerequisite pointer invoked at proof line(s) 2489; it names a standard fact of the same chapter that the counted argument uses, and every step of the counted proof built on it is retained verbatim in the delivered body.\par\medskip
\ifdefined\hypertarget\hypertarget{sourcepointer-4}{}\fi
\noindent\textit{Source reference.} Original reference 07UY: sources/stacks-resolve/pinned-original-source.tex, source lines 1584--1584. Prerequisite pointer invoked at proof line(s) 2498; it names a standard fact of the same chapter that the counted argument uses, and every step of the counted proof built on it is retained verbatim in the delivered body.\par\medskip
\begin{lemma}
\label{lemma-sections-powers-I-rational}
In Situation \ref{situation-rational} assume
$E = X_s$ is an effective Cartier divisor.
Let $\mathcal{I}$ be the ideal sheaf of $E$. Then
$H^0(X, \mathcal{I}^n) = \mathfrak m^n$ and
$H^1(X, \mathcal{I}^n) = 0$.
\end{lemma}

\begin{proof}
We have $H^0(X, \mathcal{O}_X) = A$, see discussion following
Situation \ref{situation-vanishing}. Then
$\mathfrak m \subset H^0(X, \mathcal{I}) \subset H^0(X, \mathcal{O}_X)$.
The second inclusion is not an equality as $X_s \not = \emptyset$.
Thus $H^0(X, \mathcal{I}) = \mathfrak m$.
As $\mathcal{I}^n = \mathfrak m^n\mathcal{O}_X$ our
Lemma \ref{lemma-globally-generated} shows that $H^1(X, \mathcal{I}^n) = 0$.

\medskip\noindent
Choose generators $x_1, \ldots, x_{\mu + 1}$ of $\mathfrak m$. These define
global sections of $\mathcal{I}$ which generate it. Hence
a short exact sequence
$$
0 \to \mathcal{F} \to \mathcal{O}_X^{\oplus \mu + 1} \to \mathcal{I} \to 0
$$
Then $\mathcal{F}$ is a finite locally free $\mathcal{O}_X$-module
of rank $\mu$ and $\mathcal{F} \otimes \mathcal{I}$ is globally
generated by Constructions, Lemma
\ref{constructions-lemma-globally-generated-omega-twist-1}.
Hence $\mathcal{F} \otimes \mathcal{I}^n$
is globally generated for all $n \geq 1$. Thus for $n \geq 2$ we can
consider the exact sequence
$$
0 \to \mathcal{F} \otimes \mathcal{I}^{n - 1} \to
(\mathcal{I}^{n - 1})^{\oplus \mu + 1} \to
\mathcal{I}^n \to 0
$$
Applying the long exact sequence of cohomology using that
$H^1(X, \mathcal{F} \otimes \mathcal{I}^{n - 1}) = 0$ by
Lemma \ref{lemma-globally-generated}
we obtain that every
element of $H^0(X, \mathcal{I}^n)$ is of the form $\sum x_i a_i$
for some $a_i \in H^0(X, \mathcal{I}^{n - 1})$. This shows that
$H^0(X, \mathcal{I}^n) = \mathfrak m^n$ by induction.
\end{proof}

\end{document}
